\chapter{Pagine HTML dinamiche}

\section{Introduzione}

Una \textbf{pagina web dinamica} è una pagina web che cambia in risposta all'utente o in base ai dati, senza necessariamente essere ricaricata.

\subsection{Perché sono importanti le pagine web dinamiche?}

\begin{itemize}
    \item Offrono un'esperienza utente migliorata e interattiva.
    \item Permettono personalizzazioni basate su input dell'utente o profili utente.
    \item Riducono la necessità di ricaricare la pagina, risparmiando tempo e risorse.
\end{itemize}

\section{Pagine HTML dinamiche Client-side}

\textbf{Tecnologia Primaria}: JavaScript (con librerie/framework come jQuery, React, Angular, Vue.js).

\textbf{Caratteristiche}:
\begin{itemize}
    \item Modifica del contenuto in tempo reale, anche senza interagire con il server.
    \item Esegue script nel browser dell'utente.
    \item Adatta per interazioni rapide, animazioni e manipolazione del DOM.
\end{itemize}

\textbf{Esempi}: validazione del modulo in tempo reale, animazioni, SPA (Single Page Applications).

\section{Pagine HTML dinamiche Server-side}

\textbf{Tecnologie Primarie}: PHP, Ruby on Rails, Python (Django, Flask), Node.js.

\textbf{Caratteristiche}:
\begin{itemize}
    \item Genera contenuto basato su richieste al server, spesso interagendo con un database.
    \item Adatta per operazioni che richiedono l'accesso a dati persistenti o protetti.
\end{itemize}

\textbf{Esempi}: visualizzazione di post da un blog, autenticazione utente, generazione di report personalizzati.

\section{Flusso di elaborazione}

\subsection{Client-side}

\begin{enumerate}
    \item Richiesta HTTP
    \item Risposta HTTP
    \item Elaborazione e rendering nel browser
\end{enumerate}

\subsection{Server-side}

\begin{enumerate}
    \item Richiesta HTTP
    \item Elaborazione sul server
    \item Risposta HTTP
    \item Rendering nel browser
\end{enumerate}

\section{Confronto tra Client-side e Server-side}

\subsection{Definizione}

\begin{itemize}
    \item \textbf{Client-side}: modifica il contenuto di una pagina direttamente nel browser dell'utente.
    \item \textbf{Server-side}: genera o modifica il contenuto sul server prima di inviarlo al browser dell'utente.
\end{itemize}

\subsection{Tecnologie principali}

\begin{itemize}
    \item \textbf{Client-side}: JavaScript (es.\ React, Angular, Vue.js).
    \item \textbf{Server-side}: PHP, Ruby on Rails, Python (Django, Flask), Node.js.
\end{itemize}

\subsection{Interazione con il database}

\begin{itemize}
    \item \textbf{Client-side}: generalmente mediante API o servizi web.
    \item \textbf{Server-side}: diretta, utilizzando linguaggi come SQL.
\end{itemize}

\subsection{Tempi di risposta}

\begin{itemize}
    \item \textbf{Client-side}: spesso istantaneo, poiché non richiede una nuova richiesta al server.
    \item \textbf{Server-side}: può richiedere più tempo, a seconda della richiesta al server e del tempo di elaborazione.
\end{itemize}

\subsection{Esempi di utilizzo}

\begin{itemize}
    \item \textbf{Client-side}: animazioni, validazione di form in tempo reale, SPA (Single Page Applications).
    \item \textbf{Server-side}: generazione di pagine basate su contenuti di database, autenticazione utente, elaborazione di form lato server.
\end{itemize}

\subsection{Sicurezza}

\begin{itemize}
    \item \textbf{Client-side}: dati e logica sono esposti nel codice sorgente del browser.
    \item \textbf{Server-side}: maggior controllo sulla sicurezza e sulla logica, poiché l'elaborazione avviene lato server.
\end{itemize}

\section{Casi d'uso}

\subsection{Dinamica Lato Server}

\begin{itemize}
    \item \textbf{Autenticazione e autorizzazione}: gestione sicura delle credenziali degli utenti e controllo dell'accesso.
    \item \textbf{Interazioni dirette con il database}: operazioni CRUD (Create, Read, Update, Delete) su dati persistenti.
    \item \textbf{Elaborazione di dati sensibili}: operazioni che non dovrebbero essere esposte o accessibili dal lato client.
\end{itemize}

\textbf{Vantaggi}: maggiore controllo sulla sicurezza, logica nascosta dal client, gestione centralizzata.

\subsection{Dinamica Lato Client}

\begin{itemize}
    \item \textbf{Interazioni e feedback in tempo reale}: animazioni, validazioni di form immediate, transizioni.
    \item \textbf{SPA (Single Page Applications)}: applicazioni web che non richiedono ricariche complete della pagina.
    \item \textbf{Caricamenti asincroni}: recupero e visualizzazione di dati senza ricaricare l'intera pagina (ad es.\ AJAX).
\end{itemize}

\textbf{Vantaggi}: esperienza utente fluida e reattiva, riduzione delle richieste al server, interattività avanzata.
